% !TEX root = main.tex
\subsection{Our Contributions}\label{subsection:introduction-contributions}
We answer RQ\,1, RQ\,2, and RQ\,3 affirmatively.
%The contributions follow their logical dependency:
We first establish that hidden arithmetic-circuit evaluation is achievable (RQ\,1),
then define the notion that exploits this capability (RQ\,2),
and finally construct and instantiate it (RQ\,3).

\medskip
\noindent\textbf{Efficient hidden arithmetic-circuit evaluation (RQ\,1).}
The main technical challenge is proving in zero knowledge the evaluation of the function
$P(F(M, \attr, w), G(M, \attr, w))$, concealing both the function $F$ and its computation trace.
Evaluating a hidden circuit requires concealing
not only the inputs but also the circuit structure and all intermediate computation traces,
making standard non-interactive zero-knowledge arguments of knowledge (NIZKAoK) techniques for public circuits not directly applicable.
We resolve this by adapting PLONK's arithmetization~\cite{EPRINT:GabWilCio19}:
PLONK encodes each arithmetic gate as a quadratic constraint via selector values
and enforces wire consistency across gates via a permutation.
Crucially, this encoding, which we call the \emph{PLONK description} $\plonkdescription$ of $\Circuit$
and which we refer to via PLONK's arithmetization throughout,
is a vector of length linear in the total gate count of $\Circuit$.
Treating this description as part of the \emph{witness} rather than the statement
reduces hidden-circuit satisfiability to verifying a fixed public circuit $\UniversalCircuitplonk$
of the same linear size.

Concretely, to prove $F$'s evaluation without revealing its structure,
we fix $F$'s PLONK description as part of the witness,
derive random challenges $\alpha, \beta$ from the random oracle,
and view the resulting constraint system as a fixed public verification circuit
$\UniversalCircuitplonk_F$ whose size is linear in the gate count of $F$.
A NIZKAoK on $\UniversalCircuitplonk_F$ then certifies $F$'s evaluation in zero knowledge,
as efficiently as if $F$'s circuit structure were public.
This is the first result of this kind for bounded-size arithmetic circuits over $\FF_q$
without universal circuits~\cite{STOC:Valiant76}
and without the algebraic group model~\cite{CRYPTO:FucKilLos18},
and consequently yields the first efficient privacy-preserving signature scheme
supporting hidden arithmetic-circuit evaluation without universal circuits.
See Section~\ref{section:proving-function-evaluations-via-plonk} for a detailed discussion.

Prior attempts either require non-standard assumptions~\cite{EPRINT:BonNguOzd21,PETS:CGGSS24}
or restrict to narrow circuit classes~\cite{PQCRYPTO:LNPTW21,ACISP:TNLPS24};
our technical contribution here is this application of a PLONKish-based reduction,
not general uses of zero-knowledge~proofs.

\medskip
\noindent\textbf{A new notion (RQ\,2).}
We propose \emph{Sovereign Modal Signatures} (\SMS), a notion in which both axes are governed by a single unified mechanism.
``\emph{Sovereign}'' names the per-signature autonomy of the signer,
who governs the effective policy and accountability mode at each signing,
unconstrained by a predefined menu of options.
``\emph{Modal}'' names the mode-structured accountability output,
generalizing the finite multimodal disclosure of MPS~\cite{CRYPTO:NGSY22}
to an open-ended spectrum governed by the signer's choice of $G$
and the certified structure of $F$.
``Sovereign modal'' thus names a compound property:
the signer exercises \emph{sovereign} per-signature control
over a \emph{mode}-structured policy and accountability~system.

In \SMS, each signer holds a private attribute $\attr$
and a hidden policy $F$ certified by a function-issuing authority $\FA$.
For each message $M$, the signer additionally chooses an \emph{ad hoc} function
$G \in \GFamily$, freely selected at each signing without reissuing the certificate.
$G$ is \emph{public}: any verifier can confirm that $G$ approved the message,
providing the publicly auditable policy component absent from \textsf{FS} and \textsf{PBS}.
Hiding $G$ would forfeit this auditability and double the proof cost;
the flexibility of \SMS lies in the per-signature choice of $G$, not its concealment.
In practice, $G$'s inputs or outputs may themselves carry private information,
\eg, as ciphertexts or commitments, even when $G$'s structure is public.
$G$ does not override $F$'s accountability encoding;
a public combiner $P$ composes their outputs so that both the signing policy
and the accountability mode are jointly determined:
\begin{equation}\label{equation:determining-tM-op-from-G-F-P}
(\tM, \op) = P\bigl(F(M, \attr, w),\, G(M, \attr, w)\bigr).
\end{equation}
The output $\tM$ determines what is signed,
and $\op$ determines what the authority recovers.
A special value $\op \notin \OpeningSetSignable$ blocks signing entirely;
any $\op \in \OpeningSetSignable$ permits a signature and encodes the accountability mode,
ranging from unconditional anonymity through partial disclosure to full identity recovery.
The signer produces a signature on $\tM$
in the manner of a functional signature~\cite{PKC:BoyGolIva14},
binding the output to the certified hidden policy,
with $\op$ then encrypted for the tracing authority $\TA$~\cite{CRYPTO:NGSY22}.
Tracing is triggered by $\TA$ running $\SMSOpen$ on any valid signature,
recovering exactly the information encoded in $\op$.
Throughout, $\attr$ and $F$ remain hidden from all signatures.

Crucially, even with $F$ hidden, any verifier can confirm that
the signature was produced under a policy certified by $\FA$
and a publicly known $G$, with $P$ yielding a signable output,
a form of policy compliance auditability absent from FS and PBS.
The effective policy can be updated at each signing via $G$
without reissuing the key or the certified hidden $F$.
Existing primitives arise as special cases:
PBS and {FS} correspond to $P$ acting as the identity on $F$'s output;
{ABS} and {BiAS} correspond to $P$ acting as the identity on $G$'s output;
{MPS} corresponds to $G$ selecting an index into a public function family fed to $P$.
\SMS thus subsumes {RS}, {GS}, {ABS}, {BiAS},
{MPS}, {FS}, and {PBS} as special cases,
and realizes the single-signer, function-bound slice of the functional credentials framework envisioned~\cite{iotpts/Chow16}.
%in prior work~\cite{iotpts/Chow16}.

% !TEX root = main.tex
\begin{table}[t]\footnotesize\centering
\setlength{\tabcolsep}{4pt}
\caption{%
  Comparison of \SMS with prior privacy-preserving signature schemes
}
\begin{tabular}{@{} l | c l | c l | c @{}}
\hline
\multirow{2}{*}{\textbf{Scheme}}
	& \multicolumn{2}{l|}{\textbf{Public Policy}}
	& \multicolumn{2}{l|}{\textbf{Hidden Policy}}
	& \multirow{2}{*}{\textbf{Accountability}}
	\\
\cline{2-5}
	& Scoped & Class
	& Scoped & Class
	&
	\\
\hline
RS~\cite{TCC:BenKatMor06,AC:RivShaTau01}
	& --- & ---
	& --- & ---
	& ---
	\\
GS~\cite{EC:ChaVan91,KiayiasY06}
	& --- & ---
	& --- & ---
	& Identity (ID)
	\\
PBS~\cite{PKC:BelFuc14}
	& --- & ---
	& \checkmark & Linear~\cite{ChengNW16,XuSNW22}
	& ---
	\\
FS~\cite{PKC:BoyGolIva14}
	& --- & ---
	& \checkmark & NP
	& ---
	\\
ABS~\cite{RSA:MajPraRos11}
	& \checkmark & Arithmetic~\cite{SCN:Sakai22}
	& --- & ---
	& ---
	\\
BiAS~\cite{EC:LNPY21}
	& --- & ---
	& --- & ---
	& ID / None~\cite{EC:LNPY21}
	\\
\MPS~\cite{CRYPTO:NGSY22}
	& \checkmark & Arithmetic~\cite{CRYPTO:NGSY22,ACISP:TNLPS24}
	& --- & ---
	& Finite Modal
	\\
\hline\hline
\textbf{\SMS}
	& \checkmark & Arithmetic
	& \checkmark & Arithmetic
	& Sovereign Modal
	\\
\hline
%\multicolumn{6}{@{}l}%{\emph{Accountability}:
%Binary: identity or none.
%% ; Modal: open-ended spectrum.
%}
%	\\
%\multicolumn{6}{@{}l}{%
%Identity: the tracing authority recovers the signer's complete identity.}
%	\\
\multicolumn{6}{@{}l}{\emph{Scoped}: whether this policy class determines signing eligibility.}
	\\
%\multicolumn{6}{@{}l}{Any: the notion is parameterized by an arbitrary function family.}
%	\\
\multicolumn{6}{@{}l}{Arithmetic: instantiations exist for all bounded-size arithmetic circuits over $\FF_q$.}
	\\
\multicolumn{6}{@{}l}{Linear: instantiations exist for linear functions over $\ZZ_2$.}
	\\
%\multicolumn{6}{@{}l}{Sub-Boolean: instantiations exist for a proper subset of all Boolean circuits.}
%	\\
\multicolumn{6}{@{}l}{%
Finite Modal: selection from a predefined, bounded set of disclosure modes.}\\
\multicolumn{6}{@{}l}{%
Sovereign Modal: signer-governed control over an open-ended spectrum.}\\
\end{tabular}
\vspace{0.5ex}
\label{table:comparisons-with-other-privacy-preserving-signatures}
\end{table}

Table~\ref{table:comparisons-with-other-privacy-preserving-signatures} situates \SMS
among prior privacy-preserving signatures;
\SMS is the unique one with public arithmetic functions,
hidden arithmetic functions, 
sovereign per-signature policy control,
and modal accountability simultaneously.

\medskip
\noindent\textbf{Generic construction and post-quantum instantiation (RQ\,3).}
Based on the paradigm of Nguyen~\etal~\cite{CRYPTO:NGSY22},
we build \SMS generically from four standard primitives:
a secure signature scheme $\SignatureScheme$,
a chosen-plaintext-secure (CPA-secure) public-key encryption scheme $\EncryptionScheme$,
a strongly unforgeable one-time signature scheme $\OTSScheme$,
and a NIZKAoK for arithmetic circuits
satisfying simulation-sound extractability~\cite{CRYPTO:ChaLys06},
where a valid proof always yields an extractable witness
even in the presence of simulated proofs.
The hidden-circuit evaluation of NIZKAoK
is handled by the PLONK-based reduction established in RQ\,1:
fixing $F$'s PLONK description as part of the witness
reduces the proof obligation to a fixed public verification circuit,
enabling the NIZKAoK to verify $F$'s evaluation as efficiently as that of a public circuit of the same size.
The NIZKAoK proves in zero knowledge that
the signer holds a valid $\FA$-issued certificate over $(\attr, F)$,
that $P(F(\cdot, \cdot, \cdot), G(\cdot, \cdot, \cdot))$ evaluates correctly
to a signable opening tag,
and that the tag is honestly encrypted for $\TA$ and bound to
a one-time verification key certified under the signer's attribute key.
This avoids universal circuits~\cite{STOC:Valiant76},
pairing-based assumptions, and idealized algebraic models.

Regarding our post-quantum instantiation, we employ ZKBoo \cite{USENIX:GiaMadOrl16}
under the MPC-in-the-Head (MPCitH) paradigm~\cite{STOC:IKOS07}.
ZKBoo's direct handling of arithmetic-circuit satisfiability makes it a natural fit
for the NIZKAoK slot in the generic \SMS construction,
and this instantiation replaces the abstract NIZKAoK with ZKBoo
without redesigning the signing protocol.
The lattice-based primitives over $\ZZq$ include
Lindner--Peikert (LP) public-key encryption (PKE)~\cite{RSA:LinPei11} ($\EncryptionSchemeLP$),
Jeudy--Roux-Langlois--Sanders (JRS) signatures~\cite{CRYPTO:JeuRouSan23} ($\SignatureSchemeStandardJRS$),
and Kawachi--Tanaka--Xagawa (KTX) commitments~\cite{AC:KawTanXag08}.
Since ZKBoo satisfies simulation-sound extractability with quasi-unique response~\cite{INDOCRYPT:FKMV12},
\ie, any commitment and challenge admit at most one valid response up to negligible probability,
it subsumes the anti-malleability role of $\OTSScheme$ in the generic construction.

The PLONK's arithmetization induces a two-phase witness structure.
Witnesses available before the random challenges $\alpha, \beta$ are committed first.
Witnesses determined by $\alpha, \beta$ are committed in a second phase.
This requires a targeted modification to ZKBoo's decomposition,
detailed in Section~\ref{section:overview-lattice-sms}.
The resulting signature size is quasilinear in the gate count of $F$
and linear in the number of multiplications of $G$ and $P$,
with a $\log q$ overhead arising from the binary decomposition
required to parse field elements into the binary representation
compatible with the lattice-based primitives.
This framework also supports richer instantiations via
advanced MPCitH
techniques~\cite{EC:AGHHJY23,AC:AHJMRY23,PKC:BauNof20,EC:CarCouJou23,CCS:CDGORR17,CCS:KatKolWan18,CCS:DelOrsTan21,CRYPTO:FenJouRiv22}
and vector oblivious linear evaluation in the head (VOLEitH)
techniques~\cite{CRYPTO:BBdSKORS23,ACISP:Bui25,PKC:CLYYYZ24,AC:OTX24}.

Benchmarking with concrete parameter selection is deferred to future work.
See Section~\ref{section:overview-lattice-sms} and
\ifcameraready the full version~\cite{full}
\else Appendix~\ref{appendix:lattice-sms}
\fi for details
of these building blocks.
